\documentclass{article}
\title{行動経済学の落とし穴 — 選び方が歪む10の理由}
\author{TeX 図書館}

\begin{document}

\section{なぜ正しい選択ができないのか}

行動経済学は、人間の選択を観察して研究する学問である。伝統的な経済学は、人間が常に正しく計算すると仮定してきた。この仮定を\textbf{エージェント的合理性}と呼ぶ。この章では、その仮定がどこで役に立ち、どこで崩れるかを見る。あわせて、この本の地図を提示する。

\subsection{エージェント的合理性の仮定}

経済学の教科書に登場する選び手は、条件を3つ満たす。第一に、選択肢と結果をすべて知っている。第二に、自分の好みを一貫している。第三に、計算を間違えない。

\begin{quote}
\textbf{【定義】}\ \textbf{エージェント的合理性}とは、行動主体が完全な情報と一貫した選好をもち、期待する効用を最大化するように選ぶという仮定である。市場全体の分析には強い道具だが、個人の日常の判断を説明するには手足が足りない。
\end{quote}

この仮定は便利だ。企業の価格決定や為替の動きを扱うとき、細かい迷いを無視できる。モデルの線がきれいに引ける。だから捨てられたのではない。

\subsection{現実との使い分け}

実際の人間は、3つの条件をどれも満たさない。情報は膨大で、手に入らない。好みはその場で揺れる。計算は面倒なので、近道を使う。

\begin{quote}
\textbf{【例】}\ スーパーのシャンプー棚には20種類並んでいる。全銘柄の成分と価格を比べる人はいない。大半は「いつもの」「半額の」「前回買った」で選ぶ。これは怠慢ではなく、認知資源の節約である。ただし節約の副産物として、罠にも掛かりやすい。
\end{quote}

\begin{quote}
\textbf{【注意】}\ 行動経済学は「人間は不合理だ」と宣告する学問ではない。慣れた領域、結果がすぐ返ってくる領域では、人間はよく働く。問題は、\textbf{慣れていない・結果が遅い・数字が大きい}場面で、近道が誤答を出ることだ。
\end{quote}

\subsection{この本の地図}

選び方が歪む理由は、観察されている限り10ある。以後の章はこの並びに沿う。

\begin{itemize}
\item \textbf{現状維持}: 動かないほうが楽なので、放っておく(第2章)
\item \textbf{損失回避}: 失う痛みは、得る喜びより大きい(第2章)
\item \textbf{アンカー}: 最初に見た数字が、後の判断を釘付けにする(第3章)
\item \textbf{参照点依存}: 同じ結果でも、どこを基準にするかで評価が変わる(第4章)
\item \textbf{デフォルト}: 何もしなかった結果が、勝手に選ばれる(第5章)
\item \textbf{社会的証明}: みんながやっている、を根拠にする(第6章)
\item \textbf{権威}: 名前と肩書きで判断を止める(第6章)
\item \textbf{自己制御不足}: 今の自分と未来の自分が別人(第7章)
\item \textbf{選択疲労}: 選択肢が多いほど、決断の質が落ちる(第8章)
\item \textbf{情報過多}: 情報が増えると、比較の軸が消える(第8章)
\end{itemize}

この10個は、互いに組み合わさって働く。サブスクの放置は、現状維持とデフォルトと選択疲労の合流だ。詐欺的な売り込みは、アンカーと権威と社会的証明の三段編みである。第9章で、それらを一つのチェックリストに束ねる。

\section{現状維持と損失回避}

歪みの最大の供給元は、\textbf{動かないことの安らかさ}である。何もしなければ、今の状態が続く。人間はそれを好む。

\subsection{デフォルトの力}

保険、サブスク、労働条件は、加入時の一回の判断で数年動かない。更新の時期が来ても、見直す手間は確実に、放っておく手間より大きい。

\begin{quote}
\textbf{【例】}\ 手機のキャンペーンプランは、当初だけ安く、翌月から通常料金に戻る。解約にはアプリを開き、メニューを2つ潜り、確認画面を通す必要がある。そのまま通常料金で半年使う人は多い。「決めたからではなく、決めなかったから」支払っている典型的例である。
\end{quote}

労働条件も同じだ。固定残業代の適用、勤務地の変更、評価制度の改定。会社が一方的に変更しない限り、社員は現状を維持する。異議を唱えるコストの高さが、同意の雰囲気を作る。

\subsection{損失は得の2倍重い}

得が1回の喜びなら、損は1回の痛みではない。\textbf{失う感覚は、得る感覚の約2倍}と言われる。これは比喩ではなく、行動観察から繰り返し確認された傾向である。

\begin{quote}
\textbf{【例】}\ ポイントの失効期限が近づくと、普段買わない物まで買ってしまう。「使い切らないと損」という感覚は、実際には予定になかった出費である。売り手はこの点を知っている。だから「今だけ」「残り3日」を文字面ではなく、心理的な締め切りとして使う。
\end{quote}

損失回避は投資でも出る。含み損の株を抱え続け、含み益の株を早く売る。損を確定させるのが痛いので、逆にしている。売却を先送りしている間、資金は拘束され、判断は歪んだままである。

\begin{quote}
\textbf{【注意】}\ 「得しよう」という売り文句より、「失わせる」という売り文句のほうが強い。ただし反転も可能だ。自分が誘導されたと気づいたら、得の言い方に直して読み直す。「今買うと3000円得をする」を「今買わなければ3000円損をする」と言い換え、それでも必要なら買う。
\end{quote}

\begin{quote}
\textbf{【演習】}\ 自分が現在払い続けているサブスクを3つ挙げ、それぞれについて「先月、自分で能動的に更新したか」「何もしなかった結果か」を分類せよ。
\end{quote}

\textbf{解答の指針:} 能動的に更新したものが1つもなければ、その契約は現状維持で維持されている。次に、解約にかかる手数と手間を書き出し、手間が3分以内なら今月中に検討、それ以上なら解約画面まで進む、のどちらかを先に決めておく。判断材料は金額ではなく、\textbf{更新のとき自分が行動したか}である。

\section{アンカー}

最初に見た数字は、その後の判断の土台になる。土台が傾いていれば、建物も傾く。

\subsection{相場・給与・値引き}

人間は価格を絶対値では測れない。\textbf{前の数字と比べて}測る。だから最初に提示された数字が、正当な価格の見積もりを支配する。

\begin{quote}
\textbf{【例】}\ 古着店で値札9800円のジャケットに5800円の値札が付いていたら、安く感じる。隣の店で同じ服が4000円だとしても、最初の9800円が頭に残ったまま比べることになる。相場観は「客観的な出発点」ではなく、店が置いてきた1枚の札から始まる。
\end{quote}

給与交渉も同じ構造だ。最初に求職者が提示する金額が、その後の交渉レンジを固定する。早めに言うと不利、と言う経験則は、アンカーが効くから生まれた。会社側が先にレンジを出さないのも、同じ理由である。

値引き相談では、こちらが先に正直な予算を言うと、それがアンカーになる。逆に、相手の初値を聞かずに自分の根拠から積み上げた数字を先に置くと、こちら有利に転ぶことがある。

\subsection{アンカーから逃れる技術}

アンカーは消せない。消せないなら、\textbf{別のアンカーを先に自分で置く}。

\begin{itemize}
\item 買う前に、\textbf{自分が払ってよい上限}を1行で書く。店を入る前に書くのが肝要。
\item 値札と同時に、\textbf{他店の同型番号}を1つだけ確認する。比較対象を1つ持つだけで釘は緩む。
\item 「定価」「通常価格」「相場」の3つを分けて書き、どれが自分の判断根拠かを決める。
\end{itemize}

\begin{quote}
\textbf{【注意】}\ 割引率の大きい順に店を選ぶ方法は、割引率自体がアンカーになる。「70\%オフ」は元値が正しくて初めて意味を持つが、元値は店が書いたものである。割引率ではなく、\textbf{支払額そのもの}と必要度で比べること。
\end{quote}

\begin{quote}
\textbf{【演習】}\ 中古パソコンの購入で、希望機能がすべて揃う上限価格を先に決めよ。実際に店頭で、高い機種がアンカーとして置かれた場面を想定し、自分の数字を守る具体的な行動を2つ書け。
\end{quote}

\textbf{解答の指針:} 上限価格は店に入る前にメモしておく(例: 用途から逆算した上限)。守る行動は、たとえば「店員に触られる前に自分の数字を同僚に送る」「高い機種の性能表ではなく自分の必要機能リストを開く」など、\textbf{アンカーを入れる前に自分の基準を再提示する}ものが有効。金額の妥当性を検討し直すことではなく、比較の出発点を自分に戻すこと。

\section{プロスペクト理論の勘所}

得と損は、鏡に映った対称ではない。同じ金額でも、立場によって意味が変わる。この章では数式を使わず、言葉だけで骨格を掴む。

\subsection{参照点}

評価の出発点を\textbf{参照点}と呼ぶ。給料、購入価格、前年との比較、期待していた結果。人間は結果の絶対値ではなく、参照点からの差で喜びと悲しみを測る。

\begin{quote}
\textbf{【定義】}\ \textbf{参照点}とは、得と損を判定する基準となる水準である。同じ50万円の昇給でも、同僚が80万円昇給していたら損に感じられる。結果の意味は金額ではなく、参照点との距離が決める。
\end{quote}

参照点は本人が勝手に置くので、客観的でない。だから操作できる。相手に「他社平均はこれ」を見せれば、それがあとの話の基準になる。

\subsection{得と損で変わるリスク選好}

参照点より上で(得の領域)、人間は\textbf{確実な得}を選ぶ傾向がある。確率75\%で1万円を得る案より、確実に7000円を得る案を選ぶ人が多い。

参照点より下で(損の領域)は逆になる。確実な損を避けて、\textbf{賭けに打って出ようとする}。確実に7000円を損ねる案より、確率75\%で1万円を損ねる案を選ぶ人が多い。

\begin{quote}
\textbf{【例】}\ 株で30万円の含み損を抱えた人は、危ない銘柄を抱え続ける。損を確定させる代わりに、取り返せるかもしれない賭けに望む。一方、20万円の含み益が出た人は早く売る。ここでは「賭けるのが好き」でも「慎重」でもない普通の人ほど、この非対称に動く。
\end{quote}

もう一つの勘所は\textbf{確率の扱い}である。小さな確率は誇張され、大きな確率は見劣りする。宝くじは実際の期待値より魅力的に見え、99\%の当選確率は100\%に見劣りする。「ほぼ確実」は、実際より少し弱く感じられる。

\subsection{「無料」の引力}

参照点の操作に、最強の道具が「無料」である。無料は価格の一種ではなく、\textbf{損の領域から得の領域への強制移送}として働く。0円の右側には、何も失わないという安心がある。

\begin{quote}
\textbf{【例】}\ 送料無料の条件が商品価格3000円だとすると、2980円の物を1つ足して3000円に合わせる客が出る。足した分は支払であり、送料を節約したわけではない。この上乗せは、送料を「失うもの」として枠を作ったから生まれる。無料が1円でもあれば、人は比較を止める。
\end{quote}

対策は単純だ。\textbf{「無料」を1つの数字として扱わない}。「無料で付く付属品」「無料で使える期間」「無料のあとに続く料金」の3つに分けて書き、支払総額だけで比べる。無料は条件の一種であり、保証ではない。

\begin{quote}
\textbf{【注意】}\ 「初月無料」は、無料の翌月に課金がある限り、2つの時点の比較である。初月の喜びだけを見て判断すると、失う側に立たされる。翌月の金額を初月と並べて書くだけで、判断は素直に戻る。
\end{quote}

\begin{quote}
\textbf{【注意】}\ 悪質な販売は、客を損の領域にわざと押し込む。「今契約しなければ毎月2000円損」と言葉を変えると、客は賭けに打って出ようとする。冷静さは、\textbf{参照点を自分の持ち場に戻す}ことでしか戻らない。持ち場は「今いくら得するか」ではなく「そもそも私はそれを必要としていたか」である。
\end{quote}

\section{デフォルトと誘導(ナッジ)}

何もしなければ選ばれる選択肢を\textbf{デフォルト}という。設計者は、そこに何を置くかを決められる。この権力は、良い形にも悪い形にも使われる。

\subsection{良い設計と悪意ある設計の見分け方}

良い誘導は、本人の利益と選択の自由を両立させる。悪い誘導は、手間の非対称で離脱を妨げる。

良い設計の例は多い。医療機関が健康診断の予約を済ませてから連絡する、電力の切り替えで白紙のままなら最安プランになる、社員は入社時に退職金制度の掛率選択を促される。いずれも、\textbf{望ましい行動がデフォルト}になっている。

\begin{quote}
\textbf{【例】}\ 動画配信サービスの無料体験は、開始時に支払い登録が必須で、終了後に自動で課金される。解約は3クリック、開始は1クリック。この非対称そのものが、設計者の意図の可視化である。逆に、開始と解約の手数が同じ設計は、少なくとも離脱を妨げる意図がない。
\end{quote}

見分け方は3つの問いに絞れる。

\begin{itemize}
\item \textbf{誰の利益か}: デフォルトは、本人と事業者のどちらに都合がいいのか。
\item \textbf{切り替えは対称か}: 入る手間と出る手間は同じか。違うなら、差が離脱防止の費用だ。
\item \textbf{情報は一目で見えるか}: 課金開始日、解約手順、条件変更が最初の画面にあるか。
\end{itemize}

\subsection{オプトアウトの罠}

加入を自動にし、出るには意思表示を要する方式を\textbf{オプトアウト}という。逆に、何もしなければ加入しない方式をオプトインと呼ぶ。

オプトアウトは合理的に見える。面倒な手続きを減らすからだ。だがそれは、\textbf{何もしなかった人が勝手に契約する}ことでもある。高齢者の保守サービス、解約条件の厳しい通信契約、暗黙の同意チェックボックス。いずれもこの型である。

\subsection{良い誘導の条件}

誘導そのものは悪ではない。悪いのは、\textbf{離脱だけを難しくする}設計である。良い誘導には共通点がある。

\begin{itemize}
\item デフォルトは、本人の利益に沿う選択になっている。
\item そう書かれた理由が、最初の画面に一文で添えられている。
\item 切り替えのコストは、入るコストと同じ大きさである。
\item 後から見直せる。契約の日付と解約先が、いつでも表示される。
\end{itemize}

4つ目が最も重要だ。\textbf{後から直せる設計}は、本人が自分の選択に責任を持てる。1クリックで始め、1クリックでやめられるなら、デフォルトは便利な省力化である。

\begin{quote}
\textbf{【注意】}\ 「申し込むのが簡単だから便利」という説明は、片側にしか働かない。申し込む手間と、やめる手間が違うとき、やめる手間の分だけ、不要な契約が残る。申込みの1クリックを称賛する文面は、同時に解約の5クリックを隠している。
\end{quote}

\begin{quote}
\textbf{【演習】}\ 次の3つを、オプトインとオプトアウトに分類せよ。(1)退職金の企業型DCで、何もしなければ加入する(2)無料ニュースレターで、登録した人だけ届く(3)団体保険で、希望者のみ加入する。
\end{quote}

\textbf{解答の指針:} (1)は何もしなければ加入なのでオプトアウト、(2)は登録が条件なのでオプトイン、(3)は希望表明が必要なのでオプトイン。分類の基準は\textbf{何もしなった場合の状態}である。自社や自分が関わる契約をこの基準で見直すと、放置で発生している義務が見える。

\section{社会的証明と権威}

判断に自信がないとき、人間は他人の行動と肩書きに読み込む。速いが、精度は保証されない。

\subsection{みんながやっている}

同じ状況の他人の行動は、手っ取り早い情報である。ラーメン店の行列、レビューの件数、販売数の表示。これらは\textbf{社会的証明}と呼ばれる。

\begin{quote}
\textbf{【例】}\ 通信販売の商品ページに「本日184名が購入」と出ると、クリック率が上がる。この数字は、商品の良し悪しを何も示していない。人気の理由が価格なのか品質なのか、自らの用途に合うのか、すべて不明のまま数字だけが効いている。
\end{quote}

社会的証明が効くのは、\textbf{自分に判断材料がない}ときである。効く理由が正当かを問わずに流れる点が危険だ。行列は行列であり、味の保証ではない。仕入れの少なさ、席数の少なさ、名物の書き込みやすさでもある。

\subsection{権威}

肩書き、資格、制服、専門用語。これらは判断を速くする。同時に、\textbf{検証を止める}。

\begin{quote}
\textbf{【例】}\ 「当社認定の資格取得で収入が上がる」という勧誘で、認定しているのが販売元自身なら、話は閉じている。肩書きが本物でも、\textbf{その人がこの件の利害から離れているか}は別問題である。推薦記事の末尾に広告表記がない場合、推薦は商品購入で成立している可能性が高い。
\end{quote}

権威を検証する手順は短くてよい。

\begin{itemize}
\item その肩書きは\textbf{誰が発行し、何を保証する}のか。
\item その人はこの件で\textbf{お金の利害}をもつか。
\item \textbf{独立した第三者}の同じ結論があるか。
\end{itemize}

この3つに答えられない権威は、飾りである。飾りを根拠に大金を動かさないこと。

\begin{quote}
\textbf{【注意】}\ 社会的証明と権威は、詐欺的な勧誘で必ず使われる。順番も決まっている。まず権威で信頼を作り、次にみんなが買っていると示し、最後に損失回避で締める。3つの符号が同時に見えたら、それは誘導の型であって、中立的な情報ではない。
\end{quote}

\subsection{記事と広告の境界}

信頼を作ったあとに金を取る装置が、\textbf{記事と広告の混在}である。体験談、ランキング、比較記事のどれが広告かは、見出しの小文字の表記でしか分からないことが多い。末尾に「PR」「広告」の表示があれば、その記事は売り手が用意したものだと考える。

\begin{quote}
\textbf{【例】}\ 「専門家が選ぶ3選」という記事で、1位だけ購入リンクが付き、2位と3位にリンクがない場合、選定基準は編集ではなく収益である。リンクの有無は、\textbf{どの候補からお金が得られるか}の地図として読める。並び順や推薦の言葉より、この差が本音を示す。
\end{quote}

\section{自己制御の曲がり}

「やればよいと分かっているのにやらない」は、性格の問題ではない。設計の問題である。

\subsection{今の自分と未来の自分が別人}

計画する自分は真面目だ。実行する自分は忙しい。この2人の間には、\textbf{時間による価値の目減り}がある。将来の報酬ほど、今より軽く感じる。だから運動も、貯蓄も、後回しになる。

\begin{quote}
\textbf{【例】}\ 夜スマホを床で充電し、朝アラームで起こされる人は、通知を確認するまで布団から出ない。充電場所の位置が、起床後の行動を決めている。意志の強さではなく、\textbf{スロットの配置}が問題である。
\end{quote}

\subsection{時間割・場所・仕組み}

人間の意志を補うのは、意志ではなく仕組みだ。有効な型は3つある。

\begin{itemize}
\item \textbf{時間割}: 曜日と時刻を固定する。「週2」ではなく「火曜20時にジムの前で集合」。
\item \textbf{場所}: 誘惑の物理的距離を変える。寝室に仕事道具を持ち込まない、スマホを別室に置く。
\item \textbf{仕組み}: 本人の判断を省略する。給料日の翌日に自動で別口座へ振替える先取り貯蓄。
\end{itemize}

先取り貯蓄が効く理由は、金額の賢さではない。\textbf{残高を見ずに済む}ことにある。毎月自分で判断する項目を1つ減らすと、判断の失敗も1回減る。同じ原理が、薬の定時服用、通信費の上限設定、タバコの自動引き落としにも使える。

\begin{quote}
\textbf{【例】}\ 給料日の翌日に、手取りの一定割合を別口座へ自動振替える設定は、多くの銀行アプリでできる。重要なのは割合ではなく、\textbf{実行する日を給料日に固定する}ことである。「余ったら貯める」は、毎月残高を見る判断が毎回必要な方式だ。判断の回数は、失敗の回数でもある。
\end{quote}

\begin{quote}
\textbf{【例】}\ 未来の自分との契約として、失敗したら寄付金を払うという仕組みがある。期限を決めて、未達なら指定の団体に金を渡す。支払いの痛みが、後回しの安楽さを上回るよう設計する。仕組みは強制力を持たないが、\textbf{放置の代価}を可視化する点で有効である。
\end{quote}

\begin{quote}
\textbf{【演習】}\ 自分で「今月中にやる」と言い続けている行動を1つ選び、時間割・場所・仕組みの3観点から対策を各1行で書け。
\end{quote}

\textbf{解答の指針:} 例として「読書」を選び、「火曜と金曜の21時に机で20分(時間割)」「本はリビングのテーブルに開いたまま置く(場所)」「読了記録アプリに毎晩自動リマインドを設定(仕組み)」。有効かどうかは、\textbf{実行の瞬間に意志を要するか}で測る。要しない案ほど強い。意志を要する案は、忙しい日に必ず破綻する。

\section{選択疲労と情報過多}

選択肢は、多いほど良いとは限らない。増えすぎた瞬間から、判断の質は落ちる。

\subsection{選択肢が増えるほど決断が遅い}

果物店で試した古典的な観察がある。試食の種類を6種類にすると、客は試食し、その後買った。種類を24種類に増やしたら、試食する客は増えたのに、買う客は減った。選べる喜びと、選べない停滞が同居している。

\begin{quote}
\textbf{【例】}\ ジョブサイトで条件を指定せず検索すると、5000件が並ぶ。求職者は最初の30件を眺め、疲れ、そのまま閉じる。絞り込んで12件にしたら、比較ができ、応募が決まる。\textbf{件数は情報量であり、仕事量でもある}。
\end{quote}

遅れる理由は2つある。比較の軸が見えなくなること。そして、選んだ後の後悔への備えが膨らむこと。選択肢が多ければ「他にあったかもしれない」という懐疑が残り、決断は重くなる。

\subsection{デフォルト絞り込みの技術}

大規模な選択には、\textbf{絞り込みのデフォルト}を先に置く。順番が重要だ。

\begin{itemize}
\item \textbf{譲れない条件を先に3つ決める}: 価格上限、納期、場所。この3つで機械的に落とす。
\item \textbf{比較対象を3つまでに固定}: 4つ目から先は、判断を研ぎ澄ませず疲れさせるだけ。
\item \textbf{決断の時間を先に予約する}: 火曜の19時だけを決断時間にする。それ以外の時間は選択肢を開かない。
\item \textbf{寝かせる}: 重大な契約は一晩置く。ただし「いつまでに決めるか」を先に書く。無期限の先送りは別の失敗である。
\end{itemize}

\begin{quote}
\textbf{【注意】}\ 「朝が一番良い判断ができる」という通説は当てにならない。重要なのは時間帯ではなく、\textbf{疲れと空腹を避けること}と、選択肢の数を先に減らしておくことだ。判断の質は、決める瞬間よりも、その前に何を削ったかで決まる。
\end{quote}

情報過多にも同じ処方が効く。比較サイトを開く前に、\textbf{何と何を比べるか}を3行で書く。書かずに開くと、サイト側の並び順が自分の優先順位になる。

\subsection{情報が増えると消えるもの}

情報は増えるほど便利ではない。増えると、3つが消える。

\begin{itemize}
\item \textbf{比較の軸}: 項目が増えるほど、どの項目を重く見るかが曖昧になる。
\item \textbf{締め}: 選択肢が無限なら、いつまでも決めなくてよい錯覚が生まれる。
\item \textbf{記憶}: 一度に検討した候補は、あとから思い出せなくなる。検討の記録が消えると、最初からやり直すことになる。
\end{itemize}

\begin{quote}
\textbf{【例】}\ 宅配の時間帯指定で選択肢が10個以上あると、客は待機時間を延ばしてまで空いている枠を探す。3択なら、実害のない枠をすぐ選ぶ。選択肢が減ると決断は早くなり、結果への満足は下がらない。ここでも、増やしたコストは選択側が払っている。
\end{quote}

手帳やメモへの\textbf{候補の記録}は、この歪みへの処方箋である。3つを書いたら閉じる。書いたあとの検討は、記録された3件の間でしか行わない。

\section{仕事とお金の現場で}

最後の章では、これまでの歪みを一つの運用に束ねる。給与、購入、契約の3場面で使う。

\subsection{給与交渉}

交渉はアンカーと損失回避のグラウンドである。次の順序が基本になる。

\begin{itemize}
\item \textbf{自分の数字を先に書く}: 市場相場、自分の実績、他社オファーの有無から上限と下限を2行で決める。
\item \textbf{損の言い方に備える}: 「提示額は減らされる」と言われたら、減らされる前提で可変費を下げる。
\item \textbf{現状維持を言い訳に使わない}: 前回と同じ額という会社側のデフォルトには、更新のたびに異議を唱える。
\end{itemize}

昇給面談は、\textbf{評価の振り返りと金額の交渉を分け}ると強い。同じ表でやると、評価への反論が金額の妥協に化ける。

\subsection{購入判断}

高額な買い物ほど、7つの問いで防御する。

\begin{itemize}
\item 総額はいくらか(分割の総支払、付帯費用を含む)。
\item 解約と返品の手順は、開始と同じ手間か。
\item 自動更新の日付と、それ以前にできる操作は何か。
\item 乗り換えにかかる移行コストは何か(機器、データ、違約)。
\item 自分の必要条件のうち、この商品が満たさないものは何か。
\item 3日後にまだ買う意思があるか。
\item 買わない場合、実際に何が起ころうか。
\end{itemize}

最後の2つは\textbf{損の言い方を解除する}ための問いである。買わずに困る事実が書けないなら、それは売り文句の損失であって、自分の損失ではない。

\subsection{契約書}

契約で見るべき箇所は限られる。自動更新、中途解約、違約金、条件変更の通知方法。この4つに印を付けてから、本文を読む順序を変える。\textbf{先にリスク条項}である。

\begin{quote}
\textbf{【例】}\ 「1年間ご利用無料、以後自動更新」の一文は、通常2箇所に分けて書かれる。無料期間の終わりの日付は本文中に小さく、更新の扱いは別項に書かれていることがある。検索で「更新」と「解約」を先に引くだけで、読み飛ばしの事故は減る。
\end{quote}

\subsection{詭弁的な売り文句への防御チェックリスト}

\begin{quote}
\textbf{【定義】}\ \textbf{詭弁的売り文句}とは、事実として誤でなくとも、判断の前提を偽って行動を誘導する表現である。内容ではなく、\textbf{どの歪みに効いているか}で見分ける。
\end{quote}

\begin{itemize}
\item 「今だけ」「残り2名」→ 損失回避と社会的証明。締め切りは本当に固定か。
\item 「定価から70\%オフ」→ アンカー。元値は誰が決めたか。
\item 「専門家推奨」「公式認定」→ 権威。発行者は誰か、利害はないか。
\item 「同業の9割が導入済み」→ 社会的証明。導入の理由は分かっているか。
\item 「何もしなければ現状維持です」→ 現状維持。現状は本当に維持されるか、条件は変わるか。
\item 「あとでいつでも解約できます」→ オプトアウト。解約の手順を今見せてもらえるか。
\end{itemize}

\begin{quote}
\textbf{【総合演習】}
\begin{enumerate}
\item 詭弁的売り文句6つのうち、自分の直近の購入で引っかかったものを1つ選び、効いていた歪みを特定せよ。
\item 給与交渉で相手が先に「上限はこれ」と提示してきた場合、アンカーを打ち消す行動を2つ書け。
\item 自分が契約しているサービスを1つ選び、自動更新の日付と解約手順のクリック数を実際に確認せよ。
\end{enumerate}
\end{quote}

\textbf{解答の指針:} 1. 歪みの特定は、売り手の言葉が「失う」「みんな」「専門家」「定価」のどれを指しているかで判定する。2. 有効なのは、(a)自分の根拠から積み上げた数字を先に口頭で示す、(b)市場相場や他社オファーなど、独立した比較対象を1つ持ち込む、のいずれか。相手の数字との差を語り始めると、比較の出発点は相手に移る。3. クリック数が5回を超えるなら、それは解約妨害の設計である。今すぐ解約画面まで進み、最終確認の直前で止めて、手順をメモに残す。

\end{document}
